\begin{textsample}{POS Dim 1 – 2024 – Score 18.00 – 2024/t003168.txt}  \label{ex:f1_pos_050}
When I look up at \textbf{the} sky at night, I see 100 billion stars of \textbf{the} Milky Way galaxy.
They look like lights in cabins of a giant spaceship, \textbf{The} Milky Way, sailing through space.
And I wonder if there are other passengers in those cabins.
There are 100 billion of them, comparable to \textbf{the} number of people who ever lived on Earth.
It would be arrogant to think otherwise, that we are alone, that we are unique and special, especially if you read \textbf{the} news every day.
We are not \textbf{the} pinnacle of creation.
There is room for improvement.
I ' m just a curious farm boy, and I wonder about \textbf{the} world around me.
And I hate to behave like \textbf{the} adults in \textbf{the} room, because they often pretend to know more than we actually know.
And that bothered me since I was a young kid.
And so I decided to become a scientist and answer \textbf{the} questions based on evidence, not based on prejudice, not based on \textbf{the} politics of getting \textbf{the} largest number of likes on social media.
I don ' t have a footprint on social media.
I \textbf{enjoy} nature.
Whatever it brings to our doorstep is welcome.
So let ' s just look around.
And for 70 years, we ' ve been \textbf{searching} for radio signals.
This is equivalent to staying at home and waiting for a phone call that may never come because nobody cares that we are lonely.
It may also be that others are addicted to digital screens and they live in a virtual reality, as we are at this point in time.
A much better approach is to check if there is any \textbf{object} in our backyard that may have arrived from a neighbor ' s yard.
Like a tennis ball, that may tell us that \textbf{the} neighbor plays tennis.
And we haven ' t really checked until \textbf{the} last decade. \textbf{The} first \textbf{object} to have been reported by astronomers that came from outside \textbf{the} solar system looked really weird.
It was discovered by a telescope in Hawaii.
When it passed close to Earth, it was \textbf{the} size of a football field.
What you see behind me is \textbf{the} artist ' s depiction.
It looked really weird because as it was tumbling every eight hours, \textbf{the} amount of sunlight reflected from it changed by a factor of ten, which meant that it has a very extreme shape, most likely flat, like a pancake.
And moreover, it exhibited a push away from \textbf{the} sun by some mysterious force because there was no evaporation, no cometary tail around it, no dust, no gas.
So \textbf{the} question was, what is pushing it?
And I suggested that maybe it ' s \textbf{the} reflection of sunlight, but for that, \textbf{the} \textbf{object} had to be very thin, like a sail.
And that meant that it was not produced naturally.
Maybe it ' s a \textbf{surface} layer, maybe it ' s space trash, like a plastic bag tumbling in \textbf{the} wind.
And so we go back 70 years to a question that Enrico Fermi, \textbf{the} physicist, asked at Los Alamos,"Where is everybody?"Well, this is a question that single people often ask.
But if you stay at home -- You will not find anyone.
You have to go to dating sites.
At \textbf{the} very least, you need to look through your windows for other people.
And he didn ' t seek \textbf{the} evidence.
He was just asking \textbf{the} question and kept repeating it.
And if we don ' t look for evidence, we will not find anything.
It ' s a self-fulfilling prophecy.
It ' s a way to maintain our ignorance.
And science is better than politics.
We can find \textbf{the} evidence if we allocate \textbf{the} funds for it.
This is a real image, what you see behind me, it ' s \textbf{the} Tesla Roadster car that was put as a dummy payload on \textbf{the} Falcon Heavy launch of 2018.
It ' s now moving in an elliptical orbit around \textbf{the} Sun, and perhaps in 20 million years, it will collide with Earth.
And if it will do so unexpectedly, some of my colleagues would argue"This is a \textbf{rock} of a type that we ' ve never seen before."We cannot see it with our best telescopes because it ' s too small.
It doesn ' t reflect enough sunlight. ' Oumuamua was \textbf{the} size of a football field, big enough for us to see.
And so \textbf{the} next Copernican revolution would be that we are not at \textbf{the} intellectual center of \textbf{the} \textbf{universe}.
Not only that we are not at \textbf{the} physical center of \textbf{the} \textbf{universe}, but actually, you know, we arrive to \textbf{the} play relatively late.
We are not at \textbf{the} center of stage. \textbf{The} play is not about us.
We should be modest.
We keep thinking that it ' s about us, but it ' s not.
And we better find other actors that will tell us what \textbf{the} play is about.
And people often say extraordinary claims require extraordinary evidence, but they are not seeking \textbf{the} evidence.
Actually, extraordinary evidence requires extraordinary funding.
Elon Musk argued recently,"I don ' t see any aliens."But new scientific knowledge does not fall into our lap.
We had to invest 10 billion dollars in \textbf{the} Large Hadron Collider in order to find \textbf{the} Higgs boson.
We had to invest 10 billion dollars in \textbf{the} Webb telescope in order to find \textbf{the} first generation of galaxies.
This is \textbf{the} way science is done.
You need to put \textbf{the} effort in order to find something new.
And only over \textbf{the} past decade, we discovered \textbf{objects} that came from outside \textbf{the} solar system. \textbf{The} first one was actually a decade ago.
It was a meteor, an \textbf{object} half a meter in size that collided with Earth and burned up in \textbf{the} atmosphere.
It was spotted by US government satellites. \textbf{The} fireball that it generated released a few percent of \textbf{the} Hiroshima atomic bomb energy, and it was moving too fast to be bound to \textbf{the} Sun ' s gravity.
And so we \textbf{concluded} it ' s interstellar.
It came from outside \textbf{the} solar system.
Could it be a Voyager-like meteor?
Imagine our own spacecraft colliding with a \textbf{planet} like Earth.
In \textbf{the} future, it would appear as a meteor of unusual material strength and unusual speed, which are exactly \textbf{the} properties of this meteor from 2014.
And then ' Oumuamua was discovered in 2017, and finally a comet appeared, also from interstellar space, was moving too fast.
And so my colleagues argued,"Well, this one looks familiar.
Doesn ' t it convince you that \textbf{the} others are natural in origin? \textbf{Rocks} of a type that we ' ve never seen before?"And I say, if I go down \textbf{the} street and I see a weird person, and after that I see a normal person, it doesn ' t make \textbf{the} weird person normal.
Now \textbf{the} US -- Director of National intelligence -- Avril Haines delivered three reports to \textbf{the} US Congress, talking about unidentified anomalous phenomena. \textbf{The} good news is \textbf{the} sky is not classified.
We don ' t need to wait for \textbf{the} US government to tell us what lies outside \textbf{the} solar system.
Their day job is national security.
My day job is figuring out what lies beyond \textbf{the} solar system.
And \textbf{the} sky is not classified.
We can answer \textbf{the} question ourselves.
So I ' m leading \textbf{the} Galileo project.
We built an observatory at Harvard University that monitors \textbf{the} sky 24 / 7, looking for \textbf{objects} that are not familiar, not birds, balloons, drones, \textbf{airplanes}, satellites.
So far, we monitored half a million \textbf{objects}.
Haven ' t found anything unusual yet, but we keep looking and we are using machine-learning software to figure out what we are looking at.
But \textbf{the} most exciting \textbf{endeavor} that I was involved in is an expedition to \textbf{the} Pacific Ocean near Papua New Guinea, to look for \textbf{the} materials from this meteor that I described before.
And \textbf{the} US Space Command issued a letter to NASA confirming at \textbf{the} 99. 999 percent that this \textbf{object} indeed originated from outside \textbf{the} solar system.
Based on its high speed, it was moving faster than 95 percent of all stars in \textbf{the} vicinity of \textbf{the} Sun.
It exploded in \textbf{the} lower atmosphere, about 90 kilometers away from Manus Island in Papua New Guinea, and that meant that \textbf{the} \textbf{object} had material strength tougher than even \textbf{iron} meteorites.
And so I led an expedition in June 2023.
You can see \textbf{the} team on \textbf{the} ship that was fittingly called Silver Star, and we used \textbf{the} sled with magnets on both sides to \textbf{search} for droplets left over from \textbf{the} \textbf{explosion} of this meteor.
And at \textbf{the} bottom left you see \textbf{the} filming crew of Netflix.
They are preparing a documentary about this research, and \textbf{the} director asked me,"Avi, it looks like you are running,"because I was jogging at sunrise, as I often do on land for a few miles.
And he said,"Are you running away from something or towards something?"And I said,"Both.
I ' m running away from some of my colleagues who have strong opinions without seeking evidence, and I ' m running towards a higher intelligence in interstellar space."Now we used this sled and collected magnetic particles from \textbf{the} ocean floor about a mile deep.
And then I brought them to my colleagues at Harvard University.
They look like metallic \textbf{spheres}, very distinct from \textbf{the} background of sand in which they were collected.
And my colleague at Harvard, Stein Jacobsen, is a world-renowned geochemist.
He used \textbf{the} electron microprobe, a mass spectrometer, in his laboratory. \textbf{The} person on \textbf{the} other side of me in this picture is a summer intern, Sophie Bergstrom, who found most of our molten droplets.
And so I called her \textbf{the} spheral hunter.
And most of our spherules were actually of a type familiar from \textbf{the} solar system, but about 10 percent of them looked unusual, and they had a chemical composition very different from solar system materials.
They had abundances of elements like beryllium, lanthanum, uranium, that are up to a factor of 1, 000 more than found in solar system materials.
They were not from \textbf{the} Earth, not from \textbf{the} moon, not from Mars, not from \textbf{asteroids}.
And so now \textbf{the} question arises, was this a \textbf{rock} from another star?
And of course, one possibility is that there was a natural process that produced it.
For example, most stars are \textbf{dwarf} stars, 10 percent of \textbf{the} mass of \textbf{the} sun, and they are 100 times denser than \textbf{the} sun.
And so if you bring a \textbf{planet} like \textbf{the} Earth close to them, they spaghettify \textbf{the} \textbf{planet}, make a stream of \textbf{rocks} that could be ejected at a speed similar to that of this meteor.
But it ' s also possible that this \textbf{object} was of \textbf{artificial} origin, in which case, if we look for bigger pieces of \textbf{the} \textbf{object}, we might find a gadget with buttons on it.
And I asked students in my class,"If we find such a gadget, should we press a button?"Now, some of my critics argued, maybe it ' s \textbf{coal} ash.
So we looked at 55 elements from \textbf{the} periodic table and found that \textbf{the} abundances of elements are very different from \textbf{coal} ash.
So it ' s not \textbf{coal} ash.
Others argued,"Maybe it was not a meteor, maybe it was a truck."Well, \textbf{the} data came from US government satellites.
We actually based our \textbf{search} region on \textbf{the} Department of Defense coordinates, and we went 26 times back and forth, \textbf{searching} that region.
So \textbf{the} next expedition, hopefully within \textbf{the} year, will \textbf{search} for bigger pieces of \textbf{the} \textbf{object}, maybe even \textbf{the} core of \textbf{the} \textbf{object}, because that could have a huge impact on humanity.
We all know \textbf{the} biblical story about Moses, who looked at \textbf{the} bush that was burning without being consumed, with religious awe, and that gave Moses \textbf{the} sense that there is a superhuman entity, God, out there.
Now, Friedrich Nietzsche in 1882 argued,"God is dead."And that gave rise to \textbf{the} modern period of science and technology, where humans have this hubris, they lack modesty.
Nobody is smarter than us, we are at \textbf{the} top of \textbf{the} food chain.
Maybe AI will do a little better.
But AI is just a digital mirror.
It reflects our faults.
It ' s nothing better than us.
It ' s not a digital species.
It ' s just us.
And if we find a partner out there, of course, that will give a new meaning to our existence.
And it ' s a whole different ball game.
Something from another star has nothing to do with us.
And we better be ready for that.
Not look at \textbf{the} mirror and imagine something like it, as science fiction stories do.
Now \textbf{the} good news is, next year, \textbf{the} Rubin Observatory in Chile will survey \textbf{the} southern sky every four days with a camera that is \textbf{the} size of a person, 3. 2 billion pixels, a thousand times more than your cell phone camera.
And so if we find more \textbf{objects} like ' Oumuamua, it might give us a sense of modesty.
We might bring back this sense of awe that Moses had.
Except in this time, it will be based on something that was delivered from interstellar space, from a neighbor.
And that is quite promising, actually, because it may change our priorities.
Instead of spending four trillion dollars a year on military budgets, killing each other for territories on this \textbf{rock}, \textbf{the} tiny \textbf{rock}, left over from \textbf{the} formation of \textbf{the} sun, we might realize that there is a smarter kid on \textbf{the} block, and that kid may provide a better role model than our politicians.
And if we allocate four trillion dollars a year to space exploration, we could send a CubeSat towards every star in \textbf{the} Milky Way galaxy, hundreds of billions of them, within one century.
And it gets better than that, because if we find a superhuman intelligence out there...
We might learn new physics. \textbf{The} first question I would ask is,"What happened before \textbf{the} Big Bang?"And they might have quantum gravity engineers that are capable of creating a baby \textbf{universe} in \textbf{the} laboratory.
And this job of creating a new \textbf{universe} can be perfected.
And that would help us actually, given that there is a lot of room for improvement in \textbf{the} world that we inhabit.
Thank you.
Chris Anderson : Thank you, Avi.
I ' m here like a visitor from another \textbf{planet}, because we ' re out of time.
But also just to ask you one yes or no question, OK?
If you had to put money on it, in \textbf{the} next ten years, what do you think -- is it likely that we will discover convincing evidence that persuades a critical mass of your colleagues that there is someone out there.
And I just want a yes or no.
Avi Loeb : Yes.
CA : Yes!
Avi Loeb!

% matched lemmas: airplane, artificial, asteroid, coal, conclude, dwarf, endeavor, enjoy, explosion, iron, object, planet, rock, search, sphere, surface, the, universe
\end{textsample}
